%% ECON 282E -- Session 5, Deck A3: First Order -- Four Equations, Four Unknowns
%% Source: Quant_Macro Lec.9 frames 16-24 (L461-690).
%% Exposition follows Fernandez-Villaverde & Guerron-Quintana, Solution Methods V.
%% Numbers from tools/figures/s05_numbers.json, keys "pert_first_order",
%% "higher_order", "brock_mirman".

%% ECON 282E -- Foundations of Macroeconomics (UC Riverside, Fall 2026)
%% Shared Beamer preamble for every deck in slides/S01 ... slides/S10.
%%
%% Usage (a deck lives in slides/S0X/ and is compiled from that folder):
%%   %% ECON 282E -- Foundations of Macroeconomics (UC Riverside, Fall 2026)
%% Shared Beamer preamble for every deck in slides/S01 ... slides/S10.
%%
%% Usage (a deck lives in slides/S0X/ and is compiled from that folder):
%%   %% ECON 282E -- Foundations of Macroeconomics (UC Riverside, Fall 2026)
%% Shared Beamer preamble for every deck in slides/S01 ... slides/S10.
%%
%% Usage (a deck lives in slides/S0X/ and is compiled from that folder):
%%   \input{../preamble/course_preamble.tex}
%%   \decktitle{1}{A1}{Who I Am \& Why This Course}{September 24, 2026}
%%   \begin{document}
%%   \makedecktitle
%%   ...
%%
%% Adapted from Courses/AI-ECON-2026/source/shared_preamble.tex (Metropolis theme,
%% concept/caution/example/takeaway boxes, \sectiondivider). Changes for this course:
%%   * course title block (\decktitle / \makedecktitle) instead of \makebeamertitle
%%   * \zh{...} typesets Chinese (book titles) under pdflatex via CJKutf8
%%   * researchbox for the "Research in the wild" frames
%%   * section pages off; TikZ libraries and the course map loaded here
%% Build with pdflatex only (two passes); tools/build_slides.py does this for you.

\documentclass[10pt,english,aspectratio=169]{beamer}

\usepackage[T1]{fontenc}
\usepackage{textcomp}
\usepackage[utf8]{inputenc}
\usepackage{babel}
\usepackage{CJKutf8}

\usepackage{amsmath, amssymb, amsthm, graphicx}
\usepackage{booktabs, multirow, tabularx}
\usepackage{tikz}
\usetikzlibrary{positioning,arrows.meta,shapes,calc,fit,backgrounds}
\usepackage{pgfplots}
\pgfplotsset{compat=1.16}
\usepackage[most]{tcolorbox}
\tcbuselibrary{skins, breakable}
\usepackage{listings}
\usepackage{xcolor}

\lstset{
  basicstyle=\ttfamily\small,
  keywordstyle=\color{blue!80!black}\bfseries,
  commentstyle=\color{green!50!black}\itshape,
  stringstyle=\color{orange!80!black},
  backgroundcolor=\color{gray!8},
  frame=single,
  rulecolor=\color{gray!40},
  breaklines=true,
  showstringspaces=false,
  numbers=none,
}

\usetheme{metropolis}
\usefonttheme{professionalfonts}
\metroset{sectionpage=none}

\definecolor{LightGray}{RGB}{230,230,230}
\definecolor{RedTitle}{RGB}{0,0,200}
\definecolor{VeryLightGray}{RGB}{245,245,245}
\definecolor{DarkText}{RGB}{0,0,0}
\definecolor{UCRBlue}{RGB}{0,45,106}
\definecolor{UCRGold}{RGB}{241,171,0}

% Stage colours of the course arc (used by the course map and elsewhere)
\colorlet{StageTrunk}{blue!12}
\colorlet{StageBranch}{cyan!14}
\colorlet{StageMachine}{gray!18}
\colorlet{StageWall}{red!14}
\colorlet{StageTools}{orange!20}
\colorlet{StageData}{green!16}
\colorlet{StageAgents}{violet!14}

\hypersetup{bookmarks=false,breaklinks=true,pdfborder={0 0 0},colorlinks=true,
  linkcolor=DarkText,urlcolor=RedTitle}

\setbeamercolor{background canvas}{bg=white}
\setbeamercolor{title}{fg=RedTitle}
\setbeamercolor{frametitle}{bg=VeryLightGray,fg=DarkText}
\setbeamercolor{normal text}{fg=DarkText}
\setbeamercolor{alerted text}{fg=RedTitle}
\setbeamersize{text margin left=5mm, text margin right=5mm}
\addtobeamertemplate{frametitle}{}{\vspace{-0.5em}}

\setbeamertemplate{navigation symbols}{}
\setbeamertemplate{footline}{%
  \hspace{5mm}{\usebeamerfont{page number in head/foot}\color{black!45}%
  ECON 282E \,$\cdot$\, \insertshortdeck}\hfill%
  \usebeamercolor[fg]{page number in head/foot}%
  \usebeamerfont{page number in head/foot}%
  \insertframenumber\,/\,\inserttotalframenumber\kern1em\vskip2pt%
}

% ---------------------------------------------------------------------------
% Chinese text under pdflatex (book titles etc.)
\newcommand{\zh}[1]{\begin{CJK*}{UTF8}{gbsn}#1\end{CJK*}}

% ---------------------------------------------------------------------------
% Deck title block
%   \decktitle{session}{deck id}{title}{date}
\newcommand{\insertshortdeck}{}
\newcommand{\decksession}{}
\newcommand{\deckid}{}
\newcommand{\decktitletext}{}
\newcommand{\deckdate}{}
\newcommand{\decktitle}[4]{%
  \renewcommand{\decksession}{#1}%
  \renewcommand{\deckid}{#2}%
  \renewcommand{\decktitletext}{#3}%
  \renewcommand{\deckdate}{#4}%
  \renewcommand{\insertshortdeck}{Session #1 \,$\cdot$\, Deck #1.#2}%
  \title{#3}%
  \author{Zhigang Feng}%
  \date{#4}%
}
\newcommand{\makedecktitle}{%
  \begin{frame}[plain,noframenumbering]
    \begin{tikzpicture}[remember picture,overlay]
      \fill[UCRBlue] (current page.north west) rectangle ([yshift=-0.55cm]current page.north east);
      \fill[UCRGold] ([yshift=-0.55cm]current page.north west) rectangle ([yshift=-0.62cm]current page.north east);
      \node[anchor=west, text=white, font=\small\bfseries] at ([xshift=5mm,yshift=-0.275cm]current page.north west)
        {ECON 282E \,$\cdot$\, Foundations of Macroeconomics};
      \node[anchor=east, text=white, font=\small] at ([xshift=-5mm,yshift=-0.275cm]current page.north east)
        {UC Riverside \,$\cdot$\, Fall 2026};
    \end{tikzpicture}
    \vspace*{1.1cm}
    {\usebeamerfont{title}\color{black!55}\large Session \decksession\ \,$\cdot$\, Deck \decksession.\deckid\par}
    \vspace{0.25cm}
    {\usebeamerfont{title}\color{RedTitle}\LARGE\bfseries \decktitletext\par}
    \vspace{0.7cm}
    {\large Zhigang Feng\par}
    \vspace{0.15cm}
    {\color{black!65}\deckdate\par}
    \vfill
    {\footnotesize\itshape\color{black!55}Build it on paper $\;\to\;$ solve it on the machine $\;\to\;$ push it to the frontier with AI\par}
  \end{frame}}

% ---------------------------------------------------------------------------
% Section divider (plain frame, not counted)
\newcommand{\sectiondivider}[2]{%
\begin{frame}[plain,noframenumbering]
\begin{center}
\vfill
{\Large\textbf{#1}\par}
\vspace{0.35cm}
{\normalsize #2\par}
\vfill
\end{center}
\end{frame}
}

% ---------------------------------------------------------------------------
% Boxes
\newtcolorbox{conceptbox}[1][]{colback=blue!3, colframe=RedTitle!55, boxrule=0.35mm,
  arc=1mm, left=2mm, right=2mm, top=1mm, bottom=1mm, #1}
\newtcolorbox{cautionbox}[1][]{colback=red!3, colframe=red!50!black, boxrule=0.35mm,
  arc=1mm, left=2mm, right=2mm, top=1mm, bottom=1mm, #1}
\newtcolorbox{examplebox}[1][]{colback=green!3, colframe=green!45!black, boxrule=0.35mm,
  arc=1mm, left=2mm, right=2mm, top=1mm, bottom=1mm, #1}
\newtcolorbox{takeawaybox}[1][]{colback=VeryLightGray, colframe=RedTitle!55, boxrule=0.35mm,
  arc=1mm, left=2mm, right=2mm, top=1mm, bottom=1mm, #1}
\newtcolorbox{researchbox}[1][]{enhanced, colback=violet!4, colframe=violet!60!black,
  boxrule=0.35mm, arc=1mm, left=2mm, right=2mm, top=1mm, bottom=1mm,
  fonttitle=\bfseries\small, title={Research in the wild}, #1}

\newcommand{\compacttables}{\renewcommand{\arraystretch}{1.15}}
\newcommand{\src}[1]{{\tiny\color{black!55}#1}}

% ---------------------------------------------------------------------------
\input{../preamble/course_map.tex}

%%   \decktitle{1}{A1}{Who I Am \& Why This Course}{September 24, 2026}
%%   \begin{document}
%%   \makedecktitle
%%   ...
%%
%% Adapted from Courses/AI-ECON-2026/source/shared_preamble.tex (Metropolis theme,
%% concept/caution/example/takeaway boxes, \sectiondivider). Changes for this course:
%%   * course title block (\decktitle / \makedecktitle) instead of \makebeamertitle
%%   * \zh{...} typesets Chinese (book titles) under pdflatex via CJKutf8
%%   * researchbox for the "Research in the wild" frames
%%   * section pages off; TikZ libraries and the course map loaded here
%% Build with pdflatex only (two passes); tools/build_slides.py does this for you.

\documentclass[10pt,english,aspectratio=169]{beamer}

\usepackage[T1]{fontenc}
\usepackage{textcomp}
\usepackage[utf8]{inputenc}
\usepackage{babel}
\usepackage{CJKutf8}

\usepackage{amsmath, amssymb, amsthm, graphicx}
\usepackage{booktabs, multirow, tabularx}
\usepackage{tikz}
\usetikzlibrary{positioning,arrows.meta,shapes,calc,fit,backgrounds}
\usepackage{pgfplots}
\pgfplotsset{compat=1.16}
\usepackage[most]{tcolorbox}
\tcbuselibrary{skins, breakable}
\usepackage{listings}
\usepackage{xcolor}

\lstset{
  basicstyle=\ttfamily\small,
  keywordstyle=\color{blue!80!black}\bfseries,
  commentstyle=\color{green!50!black}\itshape,
  stringstyle=\color{orange!80!black},
  backgroundcolor=\color{gray!8},
  frame=single,
  rulecolor=\color{gray!40},
  breaklines=true,
  showstringspaces=false,
  numbers=none,
}

\usetheme{metropolis}
\usefonttheme{professionalfonts}
\metroset{sectionpage=none}

\definecolor{LightGray}{RGB}{230,230,230}
\definecolor{RedTitle}{RGB}{0,0,200}
\definecolor{VeryLightGray}{RGB}{245,245,245}
\definecolor{DarkText}{RGB}{0,0,0}
\definecolor{UCRBlue}{RGB}{0,45,106}
\definecolor{UCRGold}{RGB}{241,171,0}

% Stage colours of the course arc (used by the course map and elsewhere)
\colorlet{StageTrunk}{blue!12}
\colorlet{StageBranch}{cyan!14}
\colorlet{StageMachine}{gray!18}
\colorlet{StageWall}{red!14}
\colorlet{StageTools}{orange!20}
\colorlet{StageData}{green!16}
\colorlet{StageAgents}{violet!14}

\hypersetup{bookmarks=false,breaklinks=true,pdfborder={0 0 0},colorlinks=true,
  linkcolor=DarkText,urlcolor=RedTitle}

\setbeamercolor{background canvas}{bg=white}
\setbeamercolor{title}{fg=RedTitle}
\setbeamercolor{frametitle}{bg=VeryLightGray,fg=DarkText}
\setbeamercolor{normal text}{fg=DarkText}
\setbeamercolor{alerted text}{fg=RedTitle}
\setbeamersize{text margin left=5mm, text margin right=5mm}
\addtobeamertemplate{frametitle}{}{\vspace{-0.5em}}

\setbeamertemplate{navigation symbols}{}
\setbeamertemplate{footline}{%
  \hspace{5mm}{\usebeamerfont{page number in head/foot}\color{black!45}%
  ECON 282E \,$\cdot$\, \insertshortdeck}\hfill%
  \usebeamercolor[fg]{page number in head/foot}%
  \usebeamerfont{page number in head/foot}%
  \insertframenumber\,/\,\inserttotalframenumber\kern1em\vskip2pt%
}

% ---------------------------------------------------------------------------
% Chinese text under pdflatex (book titles etc.)
\newcommand{\zh}[1]{\begin{CJK*}{UTF8}{gbsn}#1\end{CJK*}}

% ---------------------------------------------------------------------------
% Deck title block
%   \decktitle{session}{deck id}{title}{date}
\newcommand{\insertshortdeck}{}
\newcommand{\decksession}{}
\newcommand{\deckid}{}
\newcommand{\decktitletext}{}
\newcommand{\deckdate}{}
\newcommand{\decktitle}[4]{%
  \renewcommand{\decksession}{#1}%
  \renewcommand{\deckid}{#2}%
  \renewcommand{\decktitletext}{#3}%
  \renewcommand{\deckdate}{#4}%
  \renewcommand{\insertshortdeck}{Session #1 \,$\cdot$\, Deck #1.#2}%
  \title{#3}%
  \author{Zhigang Feng}%
  \date{#4}%
}
\newcommand{\makedecktitle}{%
  \begin{frame}[plain,noframenumbering]
    \begin{tikzpicture}[remember picture,overlay]
      \fill[UCRBlue] (current page.north west) rectangle ([yshift=-0.55cm]current page.north east);
      \fill[UCRGold] ([yshift=-0.55cm]current page.north west) rectangle ([yshift=-0.62cm]current page.north east);
      \node[anchor=west, text=white, font=\small\bfseries] at ([xshift=5mm,yshift=-0.275cm]current page.north west)
        {ECON 282E \,$\cdot$\, Foundations of Macroeconomics};
      \node[anchor=east, text=white, font=\small] at ([xshift=-5mm,yshift=-0.275cm]current page.north east)
        {UC Riverside \,$\cdot$\, Fall 2026};
    \end{tikzpicture}
    \vspace*{1.1cm}
    {\usebeamerfont{title}\color{black!55}\large Session \decksession\ \,$\cdot$\, Deck \decksession.\deckid\par}
    \vspace{0.25cm}
    {\usebeamerfont{title}\color{RedTitle}\LARGE\bfseries \decktitletext\par}
    \vspace{0.7cm}
    {\large Zhigang Feng\par}
    \vspace{0.15cm}
    {\color{black!65}\deckdate\par}
    \vfill
    {\footnotesize\itshape\color{black!55}Build it on paper $\;\to\;$ solve it on the machine $\;\to\;$ push it to the frontier with AI\par}
  \end{frame}}

% ---------------------------------------------------------------------------
% Section divider (plain frame, not counted)
\newcommand{\sectiondivider}[2]{%
\begin{frame}[plain,noframenumbering]
\begin{center}
\vfill
{\Large\textbf{#1}\par}
\vspace{0.35cm}
{\normalsize #2\par}
\vfill
\end{center}
\end{frame}
}

% ---------------------------------------------------------------------------
% Boxes
\newtcolorbox{conceptbox}[1][]{colback=blue!3, colframe=RedTitle!55, boxrule=0.35mm,
  arc=1mm, left=2mm, right=2mm, top=1mm, bottom=1mm, #1}
\newtcolorbox{cautionbox}[1][]{colback=red!3, colframe=red!50!black, boxrule=0.35mm,
  arc=1mm, left=2mm, right=2mm, top=1mm, bottom=1mm, #1}
\newtcolorbox{examplebox}[1][]{colback=green!3, colframe=green!45!black, boxrule=0.35mm,
  arc=1mm, left=2mm, right=2mm, top=1mm, bottom=1mm, #1}
\newtcolorbox{takeawaybox}[1][]{colback=VeryLightGray, colframe=RedTitle!55, boxrule=0.35mm,
  arc=1mm, left=2mm, right=2mm, top=1mm, bottom=1mm, #1}
\newtcolorbox{researchbox}[1][]{enhanced, colback=violet!4, colframe=violet!60!black,
  boxrule=0.35mm, arc=1mm, left=2mm, right=2mm, top=1mm, bottom=1mm,
  fonttitle=\bfseries\small, title={Research in the wild}, #1}

\newcommand{\compacttables}{\renewcommand{\arraystretch}{1.15}}
\newcommand{\src}[1]{{\tiny\color{black!55}#1}}

% ---------------------------------------------------------------------------
%% Course map: the recurring "you are here" slide for ECON 282E.
%% Loaded by course_preamble.tex. Pattern follows AI-ECON-2026 lec01_map.tex, but the map
%% shows the whole ten-session course rather than one lecture.
%%
%%   \CourseMap{s}{label}   s = 1..10 highlights session s and prints "you are here: label"
%%                          (e.g. \CourseMap{1}{1.A2}); earlier sessions get a check mark.
%%   \CourseMap{0}{}        full overview, nothing highlighted.
%%
%% Note: TikZ option lists are comma-split before TeX conditionals run, so every \ifnum
%% below sits outside [...] option lists.

\newcommand{\CourseMap}[2]{%
\begin{center}
\begin{tikzpicture}[
    sess/.style={rectangle, rounded corners=2pt, draw=black!45, minimum width=1.34cm,
        minimum height=1.12cm, text width=1.26cm, align=center, inner sep=1pt},
    stage/.style={font=\tiny\bfseries, text=black!70},
]
  \foreach \i/\d/\t/\c in {%
      1/Sep 24/Intro \\ Growth/StageTrunk,
      2/Oct 1/HA $\cdot$ OLG \\ Policy/StageBranch,
      3/Oct 8/Num.\ DP \\ Python/StageMachine,
      4/Oct 15/Numerics \\ I/StageMachine,
      5/Oct 22/Numerics \\ II/StageMachine,
      6/Oct 29/The wall \\ \textbf{Midterm}/StageWall,
      7/Nov 5/ML \\ Deep nets/StageTools,
      8/Nov 12/RL \\ HA $+$ DL/StageTools,
      9/Nov 19/Text \\ LLMs/StageData,
      10/Dec 3/Agents \\ \textbf{Projects}/StageAgents} {
    \node[sess, fill=\c] (n\i) at ({1.46*(\i-1)}, 0)
      {{\scriptsize\bfseries S\i}\\[-1pt]{\tiny\color{black!60}\d}\\[1pt]{\tiny \t}};
    \ifnum\i<#1
      \node[font=\scriptsize, text=green!45!black] at ([xshift=-2.5mm,yshift=-2.2mm]n\i.north east) {$\checkmark$};
    \fi
    \ifnum\i=#1
      \draw[RedTitle, very thick, rounded corners=2pt]
        ([xshift=-0.6mm,yshift=-0.6mm]n\i.south west) rectangle ([xshift=0.6mm,yshift=0.6mm]n\i.north east);
      % keep the label inside the picture at both ends of the row
      \ifnum\i<3
        \node[font=\scriptsize\bfseries, text=RedTitle, anchor=south west] at ([yshift=1.2mm]n\i.north west)
          {$\blacktriangledown$\,you are here: #2};
      \else\ifnum\i>8
        \node[font=\scriptsize\bfseries, text=RedTitle, anchor=south east] at ([yshift=1.2mm]n\i.north east)
          {you are here: #2\,$\blacktriangledown$};
      \else
        \node[font=\scriptsize\bfseries, text=RedTitle, anchor=south] at ([yshift=1.2mm]n\i.north)
          {you are here: #2\,$\blacktriangledown$};
      \fi\fi
    \fi
  }
  % stage band under the sessions
  \foreach \a/\b/\lab/\c in {1/1/trunk/StageTrunk, 2/2/branches/StageBranch,
      3/5/the machine $\cdot$ classical methods/StageMachine, 6/6/the wall/StageWall,
      7/8/new tools/StageTools, 9/9/new data/StageData, 10/10/colleagues/StageAgents} {
    \fill[\c] ([yshift=-1.5mm]n\a.south west) rectangle ([yshift=-4.6mm]n\b.south east);
    \path ([yshift=-1.5mm]n\a.south west) -- ([yshift=-4.6mm]n\b.south east)
      node[midway, stage] {\lab};
  }
  \node[font=\footnotesize\itshape, text=black!70] at ([yshift=-9.5mm]$(n1.south)!0.5!(n10.south)$)
    {Build it on paper $\;\to\;$ solve it on the machine $\;\to\;$ push it to the frontier with AI};
\end{tikzpicture}
\end{center}
}


%%   \decktitle{1}{A1}{Who I Am \& Why This Course}{September 24, 2026}
%%   \begin{document}
%%   \makedecktitle
%%   ...
%%
%% Adapted from Courses/AI-ECON-2026/source/shared_preamble.tex (Metropolis theme,
%% concept/caution/example/takeaway boxes, \sectiondivider). Changes for this course:
%%   * course title block (\decktitle / \makedecktitle) instead of \makebeamertitle
%%   * \zh{...} typesets Chinese (book titles) under pdflatex via CJKutf8
%%   * researchbox for the "Research in the wild" frames
%%   * section pages off; TikZ libraries and the course map loaded here
%% Build with pdflatex only (two passes); tools/build_slides.py does this for you.

\documentclass[10pt,english,aspectratio=169]{beamer}

\usepackage[T1]{fontenc}
\usepackage{textcomp}
\usepackage[utf8]{inputenc}
\usepackage{babel}
\usepackage{CJKutf8}

\usepackage{amsmath, amssymb, amsthm, graphicx}
\usepackage{booktabs, multirow, tabularx}
\usepackage{tikz}
\usetikzlibrary{positioning,arrows.meta,shapes,calc,fit,backgrounds}
\usepackage{pgfplots}
\pgfplotsset{compat=1.16}
\usepackage[most]{tcolorbox}
\tcbuselibrary{skins, breakable}
\usepackage{listings}
\usepackage{xcolor}

\lstset{
  basicstyle=\ttfamily\small,
  keywordstyle=\color{blue!80!black}\bfseries,
  commentstyle=\color{green!50!black}\itshape,
  stringstyle=\color{orange!80!black},
  backgroundcolor=\color{gray!8},
  frame=single,
  rulecolor=\color{gray!40},
  breaklines=true,
  showstringspaces=false,
  numbers=none,
}

\usetheme{metropolis}
\usefonttheme{professionalfonts}
\metroset{sectionpage=none}

\definecolor{LightGray}{RGB}{230,230,230}
\definecolor{RedTitle}{RGB}{0,0,200}
\definecolor{VeryLightGray}{RGB}{245,245,245}
\definecolor{DarkText}{RGB}{0,0,0}
\definecolor{UCRBlue}{RGB}{0,45,106}
\definecolor{UCRGold}{RGB}{241,171,0}

% Stage colours of the course arc (used by the course map and elsewhere)
\colorlet{StageTrunk}{blue!12}
\colorlet{StageBranch}{cyan!14}
\colorlet{StageMachine}{gray!18}
\colorlet{StageWall}{red!14}
\colorlet{StageTools}{orange!20}
\colorlet{StageData}{green!16}
\colorlet{StageAgents}{violet!14}

\hypersetup{bookmarks=false,breaklinks=true,pdfborder={0 0 0},colorlinks=true,
  linkcolor=DarkText,urlcolor=RedTitle}

\setbeamercolor{background canvas}{bg=white}
\setbeamercolor{title}{fg=RedTitle}
\setbeamercolor{frametitle}{bg=VeryLightGray,fg=DarkText}
\setbeamercolor{normal text}{fg=DarkText}
\setbeamercolor{alerted text}{fg=RedTitle}
\setbeamersize{text margin left=5mm, text margin right=5mm}
\addtobeamertemplate{frametitle}{}{\vspace{-0.5em}}

\setbeamertemplate{navigation symbols}{}
\setbeamertemplate{footline}{%
  \hspace{5mm}{\usebeamerfont{page number in head/foot}\color{black!45}%
  ECON 282E \,$\cdot$\, \insertshortdeck}\hfill%
  \usebeamercolor[fg]{page number in head/foot}%
  \usebeamerfont{page number in head/foot}%
  \insertframenumber\,/\,\inserttotalframenumber\kern1em\vskip2pt%
}

% ---------------------------------------------------------------------------
% Chinese text under pdflatex (book titles etc.)
\newcommand{\zh}[1]{\begin{CJK*}{UTF8}{gbsn}#1\end{CJK*}}

% ---------------------------------------------------------------------------
% Deck title block
%   \decktitle{session}{deck id}{title}{date}
\newcommand{\insertshortdeck}{}
\newcommand{\decksession}{}
\newcommand{\deckid}{}
\newcommand{\decktitletext}{}
\newcommand{\deckdate}{}
\newcommand{\decktitle}[4]{%
  \renewcommand{\decksession}{#1}%
  \renewcommand{\deckid}{#2}%
  \renewcommand{\decktitletext}{#3}%
  \renewcommand{\deckdate}{#4}%
  \renewcommand{\insertshortdeck}{Session #1 \,$\cdot$\, Deck #1.#2}%
  \title{#3}%
  \author{Zhigang Feng}%
  \date{#4}%
}
\newcommand{\makedecktitle}{%
  \begin{frame}[plain,noframenumbering]
    \begin{tikzpicture}[remember picture,overlay]
      \fill[UCRBlue] (current page.north west) rectangle ([yshift=-0.55cm]current page.north east);
      \fill[UCRGold] ([yshift=-0.55cm]current page.north west) rectangle ([yshift=-0.62cm]current page.north east);
      \node[anchor=west, text=white, font=\small\bfseries] at ([xshift=5mm,yshift=-0.275cm]current page.north west)
        {ECON 282E \,$\cdot$\, Foundations of Macroeconomics};
      \node[anchor=east, text=white, font=\small] at ([xshift=-5mm,yshift=-0.275cm]current page.north east)
        {UC Riverside \,$\cdot$\, Fall 2026};
    \end{tikzpicture}
    \vspace*{1.1cm}
    {\usebeamerfont{title}\color{black!55}\large Session \decksession\ \,$\cdot$\, Deck \decksession.\deckid\par}
    \vspace{0.25cm}
    {\usebeamerfont{title}\color{RedTitle}\LARGE\bfseries \decktitletext\par}
    \vspace{0.7cm}
    {\large Zhigang Feng\par}
    \vspace{0.15cm}
    {\color{black!65}\deckdate\par}
    \vfill
    {\footnotesize\itshape\color{black!55}Build it on paper $\;\to\;$ solve it on the machine $\;\to\;$ push it to the frontier with AI\par}
  \end{frame}}

% ---------------------------------------------------------------------------
% Section divider (plain frame, not counted)
\newcommand{\sectiondivider}[2]{%
\begin{frame}[plain,noframenumbering]
\begin{center}
\vfill
{\Large\textbf{#1}\par}
\vspace{0.35cm}
{\normalsize #2\par}
\vfill
\end{center}
\end{frame}
}

% ---------------------------------------------------------------------------
% Boxes
\newtcolorbox{conceptbox}[1][]{colback=blue!3, colframe=RedTitle!55, boxrule=0.35mm,
  arc=1mm, left=2mm, right=2mm, top=1mm, bottom=1mm, #1}
\newtcolorbox{cautionbox}[1][]{colback=red!3, colframe=red!50!black, boxrule=0.35mm,
  arc=1mm, left=2mm, right=2mm, top=1mm, bottom=1mm, #1}
\newtcolorbox{examplebox}[1][]{colback=green!3, colframe=green!45!black, boxrule=0.35mm,
  arc=1mm, left=2mm, right=2mm, top=1mm, bottom=1mm, #1}
\newtcolorbox{takeawaybox}[1][]{colback=VeryLightGray, colframe=RedTitle!55, boxrule=0.35mm,
  arc=1mm, left=2mm, right=2mm, top=1mm, bottom=1mm, #1}
\newtcolorbox{researchbox}[1][]{enhanced, colback=violet!4, colframe=violet!60!black,
  boxrule=0.35mm, arc=1mm, left=2mm, right=2mm, top=1mm, bottom=1mm,
  fonttitle=\bfseries\small, title={Research in the wild}, #1}

\newcommand{\compacttables}{\renewcommand{\arraystretch}{1.15}}
\newcommand{\src}[1]{{\tiny\color{black!55}#1}}

% ---------------------------------------------------------------------------
%% Course map: the recurring "you are here" slide for ECON 282E.
%% Loaded by course_preamble.tex. Pattern follows AI-ECON-2026 lec01_map.tex, but the map
%% shows the whole ten-session course rather than one lecture.
%%
%%   \CourseMap{s}{label}   s = 1..10 highlights session s and prints "you are here: label"
%%                          (e.g. \CourseMap{1}{1.A2}); earlier sessions get a check mark.
%%   \CourseMap{0}{}        full overview, nothing highlighted.
%%
%% Note: TikZ option lists are comma-split before TeX conditionals run, so every \ifnum
%% below sits outside [...] option lists.

\newcommand{\CourseMap}[2]{%
\begin{center}
\begin{tikzpicture}[
    sess/.style={rectangle, rounded corners=2pt, draw=black!45, minimum width=1.34cm,
        minimum height=1.12cm, text width=1.26cm, align=center, inner sep=1pt},
    stage/.style={font=\tiny\bfseries, text=black!70},
]
  \foreach \i/\d/\t/\c in {%
      1/Sep 24/Intro \\ Growth/StageTrunk,
      2/Oct 1/HA $\cdot$ OLG \\ Policy/StageBranch,
      3/Oct 8/Num.\ DP \\ Python/StageMachine,
      4/Oct 15/Numerics \\ I/StageMachine,
      5/Oct 22/Numerics \\ II/StageMachine,
      6/Oct 29/The wall \\ \textbf{Midterm}/StageWall,
      7/Nov 5/ML \\ Deep nets/StageTools,
      8/Nov 12/RL \\ HA $+$ DL/StageTools,
      9/Nov 19/Text \\ LLMs/StageData,
      10/Dec 3/Agents \\ \textbf{Projects}/StageAgents} {
    \node[sess, fill=\c] (n\i) at ({1.46*(\i-1)}, 0)
      {{\scriptsize\bfseries S\i}\\[-1pt]{\tiny\color{black!60}\d}\\[1pt]{\tiny \t}};
    \ifnum\i<#1
      \node[font=\scriptsize, text=green!45!black] at ([xshift=-2.5mm,yshift=-2.2mm]n\i.north east) {$\checkmark$};
    \fi
    \ifnum\i=#1
      \draw[RedTitle, very thick, rounded corners=2pt]
        ([xshift=-0.6mm,yshift=-0.6mm]n\i.south west) rectangle ([xshift=0.6mm,yshift=0.6mm]n\i.north east);
      % keep the label inside the picture at both ends of the row
      \ifnum\i<3
        \node[font=\scriptsize\bfseries, text=RedTitle, anchor=south west] at ([yshift=1.2mm]n\i.north west)
          {$\blacktriangledown$\,you are here: #2};
      \else\ifnum\i>8
        \node[font=\scriptsize\bfseries, text=RedTitle, anchor=south east] at ([yshift=1.2mm]n\i.north east)
          {you are here: #2\,$\blacktriangledown$};
      \else
        \node[font=\scriptsize\bfseries, text=RedTitle, anchor=south] at ([yshift=1.2mm]n\i.north)
          {you are here: #2\,$\blacktriangledown$};
      \fi\fi
    \fi
  }
  % stage band under the sessions
  \foreach \a/\b/\lab/\c in {1/1/trunk/StageTrunk, 2/2/branches/StageBranch,
      3/5/the machine $\cdot$ classical methods/StageMachine, 6/6/the wall/StageWall,
      7/8/new tools/StageTools, 9/9/new data/StageData, 10/10/colleagues/StageAgents} {
    \fill[\c] ([yshift=-1.5mm]n\a.south west) rectangle ([yshift=-4.6mm]n\b.south east);
    \path ([yshift=-1.5mm]n\a.south west) -- ([yshift=-4.6mm]n\b.south east)
      node[midway, stage] {\lab};
  }
  \node[font=\footnotesize\itshape, text=black!70] at ([yshift=-9.5mm]$(n1.south)!0.5!(n10.south)$)
    {Build it on paper $\;\to\;$ solve it on the machine $\;\to\;$ push it to the frontier with AI};
\end{tikzpicture}
\end{center}
}


\graphicspath{{./figures/}}

\lstset{language=Python, basicstyle=\ttfamily\scriptsize, frame=none,
        backgroundcolor=\color{gray!8}, aboveskip=2pt, belowskip=2pt}

\decktitle{5}{A3}{First Order: Four Equations, Four Unknowns}{Thursday, October 22, 2026}

\begin{document}

\makedecktitle

%=============================================================================
\begin{frame}{Where We Are}
\CourseMap{5}{5.A3}
\vspace{-1mm}
\begin{center}
\small \textbf{Block A:} 5.A1 what perturbation computes $\;\cdot\;$ 5.A2 the model and the parameter $\;\cdot\;$ \textbf{5.A3} first order $\;\cdot\;$ 5.A4 linear RE systems $\;\cdot\;$ 5.A5 higher order, risk, pruning $\;\cdot\;$ 5.A6 the worked example.
\end{center}
\end{frame}

%=============================================================================
\begin{frame}{Packaging Everything Into One Object}
\small
\vspace{-1mm}
\[
F(k_t,z_t;\chi)=
\begin{pmatrix}
\dfrac{1}{c(k_t,z_t;\chi)}
\;-\;\beta\,\dfrac{\alpha\,e^{\rho z_t+\chi\sigma\varepsilon_{t+1}}\big(k(k_t,z_t;\chi)\big)^{\alpha-1}}
{c\Big(k(k_t,z_t;\chi),\,\rho z_t+\chi\sigma\varepsilon_{t+1};\chi\Big)}
\\[4mm]
c(k_t,z_t;\chi)+k(k_t,z_t;\chi)-e^{z_t}k_t^{\alpha}
\end{pmatrix}
=\begin{pmatrix}0\\[1mm]0\end{pmatrix}
\]
\medskip
\begin{columns}[T]
\begin{column}{0.49\textwidth}
\begin{conceptbox}[title=\textbf{Why bother}]\footnotesize
One object, two components --- Euler equation and resource constraint --- so that ``differentiate and match'' becomes mechanical rather than a case-by-case argument.
\end{conceptbox}
\end{column}
\begin{column}{0.47\textwidth}
\begin{examplebox}\footnotesize
Write $H_i$ for the partial derivative of the Euler component with respect to its $i$-th argument. The evaluation point $(k^{*},0;0)$ is dropped when it is clear.
\end{examplebox}
\end{column}
\end{columns}
\vspace{1mm}
\src{Quant\_Macro Lec.~9, ``Defining $F(k_t,z_t;\lambda)=0$ (1/2)''.}
\end{frame}

%=============================================================================
\begin{frame}{An Identity, Not an Equation}
\begin{columns}[T]
\begin{column}{0.55\textwidth}
\small
Once the \emph{true} policies are substituted in, $F(k_t,z_t;\chi)=0$ holds for \textbf{every} $(k_t,z_t,\chi)$ in a neighbourhood. It is an identity.
\medskip

And a function that is identically zero has identically zero derivatives:
\[
  F\equiv 0 \;\Longrightarrow\; F_k=F_z=F_\chi=F_{kk}=\cdots=0 .
\]
\medskip
The unknown Taylor coefficients $\{c_k,c_z,k_k,k_z,\dots\}$ appear \emph{inside} those derivatives, through the chain rule. Setting each derivative to zero therefore produces \textbf{equations for them}.
\medskip

First order needs $(F_k,F_z,F_\chi)$. Second order also needs $(F_{kk},F_{kz},F_{zz},\dots)$, and so on.
\end{column}
\begin{column}{0.42\textwidth}
\begin{takeawaybox}\footnotesize
This single observation --- \emph{identity, therefore derivatives vanish} --- is where every perturbation equation in this course comes from.
\end{takeawaybox}
\medskip
\begin{cautionbox}\footnotesize
It holds only where the true policy is differentiable. A binding constraint breaks the identity's smoothness, not the logic.
\end{cautionbox}
\end{column}
\end{columns}
\vspace{1mm}
\src{Quant\_Macro Lec.~9, ``Why Derivatives Pin Down Unknowns'' and ``First-Order Derivatives of $F$''.}
\end{frame}

%=============================================================================
\begin{frame}{Why Doesn't $F=0$ Itself Give Us the Equations?}
\begin{columns}[T]
\begin{column}{0.55\textwidth}
\small
A fair question: $F=0$ has two components. Why are those not two equations for the first-order coefficients?
\medskip

\textbf{Because $F=0$ is the zeroth-order condition.} At the expansion point $\hat{k}_t=0$, $z_t=0$, $\chi=0$, so every linear term in the policy expansion --- $c_k\hat{k}_t$, $c_z z_t$, $c_\chi\chi$ --- multiplies zero and drops out. What survives pins down the \textbf{steady-state levels}, which we already knew.
\medskip

The slopes come from matching the coefficients on $\hat{k}_t$, $z_t$ and $\chi$, which is the same thing as imposing
\[
  F_k=0,\qquad F_z=0,\qquad F_\chi=0 .
\]
\end{column}
\begin{column}{0.42\textwidth}
\begin{conceptbox}[title=\textbf{Contrast with projection}]\small
Block B does the opposite: it imposes $F=0$ at \textbf{many points} $(k_i,z_i)$, so the unknown coefficients appear directly in the equations and no derivative is ever taken.
\medskip
One method differentiates once and evaluates nowhere; the other evaluates everywhere and differentiates nothing.
\end{conceptbox}
\end{column}
\end{columns}
\vspace{1mm}
\src{Quant\_Macro Lec.~9, ``Why Don't We Count $F=0$ Itself as Equations?''.}
\end{frame}

%=============================================================================
\begin{frame}{The First-Order Conditions}
\small
Six unknowns enter at linear order: $\{c_k,\,c_z,\,c_\chi,\,k_k,\,k_z,\,k_\chi\}$. Differentiating $F=0$ and applying the chain rule gives, schematically,
\begin{align*}
F_{k}(k^{*},0;0) &= H_{1}c_{k}+H_{2}c_{k}k_{k}+H_{3}+H_{4}k_{k}=0\\
F_{z}(k^{*},0;0) &= H_{1}c_{z}+H_{2}\big(c_{k}k_{z}+c_{z}\rho\big)+H_{4}k_{z}+H_{5}=0\\
F_{\chi}(k^{*},0;0) &= H_{1}c_{\chi}+H_{2}\big(c_{k}k_{\chi}+c_{\chi}\big)+H_{4}k_{\chi}+H_{6}=0
\end{align*}
\medskip
\begin{columns}[T]
\begin{column}{0.49\textwidth}
\begin{conceptbox}[title=\textbf{Each line is a 2-vector}]\footnotesize
$F$ has two components, so $F_k=0$ is two scalar equations, and so is $F_z=0$.
\end{conceptbox}
\end{column}
\begin{column}{0.47\textwidth}
\begin{cautionbox}\footnotesize
Note $F_k$ is \textbf{quadratic} in the unknowns: the term $H_2 c_k k_k$ is a product of two of them. First order is a linear approximation, not a linear system.
\end{cautionbox}
\end{column}
\end{columns}
\vspace{1mm}
\src{Quant\_Macro Lec.~9, ``First-order approximation''. The quadratic term is why 5.A4 needs a matrix quadratic and a root selection.}
\end{frame}

%=============================================================================
\begin{frame}{Six Unknowns, but Two Disappear}
\begin{columns}[T]
\begin{column}{0.54\textwidth}
\small
Count the restrictions. Differentiating with respect to the two \emph{state} variables gives
\[
  F_k=0 \;\Rightarrow\; 2 \text{ equations},\qquad
  F_z=0 \;\Rightarrow\; 2 \text{ equations},
\]
so four scalar equations for $\{c_k,c_z,k_k,k_z\}$ --- and nothing yet for $c_\chi,k_\chi$.
\medskip

Those two are handled separately, by
\[
  F_{\chi}=H_{1}c_{\chi}+H_{2}\big(c_{k}k_{\chi}+c_{\chi}\big)+H_{4}k_{\chi}+H_{6}=0 .
\]
This is \textbf{linear} in $(c_\chi,k_\chi)$. And if it is also \textbf{homogeneous} --- if $H_6=0$ --- then the only solution is
\[
  c_{\chi}=k_{\chi}=0 .
\]
\end{column}
\begin{column}{0.43\textwidth}
\begin{takeawaybox}[title=\textbf{Certainty equivalence}]\small
$c_\chi=k_\chi=0$ means the shock scale enters \textbf{no first-order coefficient}. To a linear approximation, an agent facing risk behaves exactly like one who does not.
\end{takeawaybox}
\medskip
\begin{examplebox}\footnotesize
So there is no precautionary saving at first order --- not because the model lacks it, but because the approximation cannot see it.
\end{examplebox}
\end{column}
\end{columns}
\vspace{1mm}
\src{Quant\_Macro Lec.~9, ``First-Order Approximation: Six Unknowns, but Two Disappear''.}
\end{frame}

%=============================================================================
\begin{frame}{Why $H_6=0$}
\begin{columns}[T]
\begin{column}{0.54\textwidth}
\small
$H_6$ is whatever sits in $F_\chi$ without multiplying $c_\chi$ or $k_\chi$ --- the constant term. It is the direct effect of $\chi$ on the equilibrium conditions, holding the policies fixed.
\medskip

At the expansion point $\chi=0$, and the innovation enters as $\chi\sigma\varepsilon_{t+1}$ inside an expectation with
\[
  \mathbb{E}_t[\varepsilon_{t+1}]=0 .
\]
A \emph{linear} change in the shock scale therefore has \emph{no} first-order effect: the term drops out.
\medskip

With no constant, the equation in $(c_\chi,k_\chi)$ is homogeneous, and the unique non-degenerate solution is zero.
\end{column}
\begin{column}{0.43\textwidth}
\begin{conceptbox}[title=\textbf{Where risk survives}]\small
Risk enters through the \textbf{variance}, which is quadratic in $\chi$. So it first appears in $c_{\chi\chi}$ --- a second-order term. 5.A5.
\end{conceptbox}
\medskip
\begin{cautionbox}\footnotesize
$\mathbb{E}[\varepsilon]=0$ is doing the work. Change the shock to something with a non-zero mean effect at first order and $H_6\neq0$.
\end{cautionbox}
\end{column}
\end{columns}
\vspace{1mm}
\src{Quant\_Macro Lec.~9, ``Why $H_6=0$ and the System is Homogeneous''.}
\end{frame}

%=============================================================================
\begin{frame}{Exactly Four Equations for Exactly Four Unknowns}
\begin{columns}[T]
\begin{column}{0.53\textwidth}
\small
Collecting:
\medskip
\begin{center}\footnotesize
\begin{tabularx}{\linewidth}{@{}>{\raggedright\arraybackslash}X r@{}}
\toprule
\textbf{Source} & \textbf{Equations} \\
\midrule
$F_k=0$, both components & 2 \\
$F_z=0$, both components & 2 \\
$F_\chi=0$ $\Rightarrow$ $c_\chi=k_\chi=0$ & --- \\
\midrule
\textbf{Total} & \textbf{4} \\
\midrule
Unknowns $c_k,c_z,k_k,k_z$ & \textbf{4} \\
\bottomrule
\end{tabularx}
\end{center}
\medskip
Square, and generically uniquely solvable. Higher-order coefficients need higher derivatives of $F$, and each order is linear given the ones below.
\end{column}
\begin{column}{0.44\textwidth}
\begin{cautionbox}[title=\textbf{``Generically''}]\small
The counting says the system is square. It does not say a \emph{stable} solution exists, or that only one does --- the system is quadratic, so it can have two roots, and only one of them is economically admissible.
\end{cautionbox}
\medskip
\begin{examplebox}\footnotesize
Choosing between them is what Blanchard--Kahn does, and it is the only genuinely new idea in Block A.
\end{examplebox}
\end{column}
\end{columns}
\vspace{1mm}
\src{Quant\_Macro Lec.~9, ``Summary: Why Exactly Four Equations for Four Unknowns''.}
\end{frame}

%=============================================================================
\begin{frame}{Measured: The Four Numbers, Computed Two Ways}
\begin{columns}[T]
\begin{column}{0.5\textwidth}
\small
On the quarterly calibration --- $\alpha=0.33$, $\beta=0.99$, $\delta=0.025$, $\rho=0.95$ --- with $k^{*}=28.3484$, the four coefficients in log deviations are:
\medskip
\begin{center}\footnotesize
\begin{tabularx}{\linewidth}{@{}>{\raggedright\arraybackslash}X r@{}}
\toprule
$\hat{k}'$ on $\hat{k}$ & $0.962061$ \\
$\hat{k}'$ on $z$ & $0.080097$ \\
$\hat{c}$ on $\hat{k}$ & $0.590408$ \\
$\hat{c}$ on $z$ & $0.322850$ \\
\bottomrule
\end{tabularx}
\end{center}
\medskip
Read the first one: capital is \textbf{very} persistent, $0.962$ per quarter, so a shock takes years to unwind.
\end{column}
\begin{column}{0.47\textwidth}
\begin{takeawaybox}[title=\textbf{Two routes, one answer}]\small
Computed by generalized Schur on the linearized system (5.A4), and independently by root-finding on $F=0$, $F_k=0$, $F_z=0$ with numerical derivatives. The two agree to $\mathbf{4.4\times10^{-11}}$.
\end{takeawaybox}
\medskip
\begin{examplebox}\footnotesize
Two implementations of the same object, disagreeing nowhere. That is the check this course keeps asking for, and it is cheap here.
\end{examplebox}
\end{column}
\end{columns}
\vspace{1mm}
\src{Ledger \texttt{pert\_first\_order}: \texttt{qz\_vs\_matching\_max\_abs\_diff} $=4.36\times10^{-11}$.}
\end{frame}

%=============================================================================
\begin{frame}{Measured: Certainty Equivalence, and What It Hides}
\begin{columns}[T]
\begin{column}{0.52\textwidth}
\small
$c_\chi$ and $k_\chi$ are the $\chi$-derivatives of 5.A2's expansion --- how much the policy shifts
when the shock is switched on. Solved at $\sigma=0$ they come out at
\[
  c_\chi = -2.7\times10^{-14},\qquad k_\chi = 3.4\times10^{-16},
\]
i.e.\ zero. So far, so predicted.

The sharper test: add the \emph{whole second-order block} and re-solve with risk on. If certainty equivalence is real, the four first-order coefficients should not move.

They move by $\mathbf{1.5\times10^{-5}}$ --- against a coefficient of $0.59$, that is one part in
$40{,}000$.

What does move is the \textbf{constant}: consumption shifts by $-0.00041\%$ of steady state. Small --- but it is the entire economics of precaution, and first order reports it as exactly zero.
\end{column}
\begin{column}{0.45\textwidth}
\begin{conceptbox}[title=\textbf{In practice}]\small
A first-order solution describes \emph{dynamics} well and \emph{risk} not at all: welfare costs of uncertainty, risk premia and precautionary saving are identically zero in it.
\end{conceptbox}
\medskip
\begin{cautionbox}\footnotesize
Asset-pricing papers that report a zero equity premium from a linearized model are not finding an economic result.
\end{cautionbox}
\end{column}
\end{columns}
\vspace{1mm}
\src{Ledger \texttt{pert\_first\_order.certainty\_equivalence} and \texttt{higher\_order.first\_order\_coeff\_shift}.}
\end{frame}

%=============================================================================
\begin{frame}{Why We Still Need Blanchard--Kahn}
\begin{columns}[T]
\begin{column}{0.55\textwidth}
\small
The first-order solution is a pair of linear rules,
\[
  \hat{c}_t=\varphi_k\hat{k}_t+\varphi_z z_t,
  \qquad
  \hat{k}_{t+1}=\psi_k\hat{k}_t+\psi_z z_t ,
\]
and substituting them into the equilibrium conditions leaves a \textbf{linear rational-expectations system}.
\medskip

Such a system can have
\begin{itemize}\small
  \item no stable solution,
  \item many stable solutions (indeterminacy),
  \item or exactly one.
\end{itemize}
\medskip
Nothing in the counting argument distinguishes these. The quadratic system of the previous slides has more than one root, and only one of them gives bounded paths.
\end{column}
\begin{column}{0.42\textwidth}
\begin{takeawaybox}\footnotesize
The coefficient count tells you the system is square. \textbf{Blanchard--Kahn tells you which root to keep} --- and whether keeping one is possible at all.
\end{takeawaybox}
\medskip
\begin{examplebox}\footnotesize
1.B1 already used this logic informally: one stable and one unstable root, so from any $k_0$ there is a unique $c_0$ on the saddle path.
\end{examplebox}
\end{column}
\end{columns}
\vspace{1mm}
\src{Quant\_Macro Lec.~9, ``Why We Need Blanchard--Kahn After First-Order Perturbation''.}
\end{frame}

%=============================================================================
\begin{frame}{Takeaways}
\begin{takeawaybox}
\begin{itemize}\small
  \item Package the equilibrium conditions as $F(k,z;\chi)=0$. Because the true policies make it an \textbf{identity}, every derivative of $F$ is zero --- and those are the equations.
  \item $F=0$ itself is the \emph{zeroth}-order condition: it pins the steady state, not the slopes. Projection (Block B) uses $F=0$ instead, at many points.
  \item Six first-order unknowns, four equations from $(F_k,F_z)$, and $F_\chi=0$ forcing $c_\chi=k_\chi=0$ because $\mathbb{E}[\varepsilon]=0$ kills the constant $H_6$.
  \item That is \textbf{certainty equivalence}: risk enters no first-order coefficient, so precaution, risk premia and welfare costs of uncertainty are all exactly zero.
  \item Measured on the quarterly model: $\hat{k}'$ loads $\mathbf{0.962}$ on $\hat{k}$ and $0.080$ on $z$; two independent solution routes agree to $\mathbf{4.4\times10^{-11}}$.
  \item Adding the whole second-order block moves those four coefficients by $1.5\times10^{-5}$ and shifts the constant by $-0.00041\%$ --- small, and the only place the economics of risk lives.
\end{itemize}
\end{takeawaybox}
\end{frame}

%=============================================================================
\begin{frame}{Bridge: Which Root, and Is There One?}
\begin{columns}[T]
\begin{column}{0.53\textwidth}
\small
We have a square system with a quadratic term, so it has more than one solution, and only one of them is economically meaningful. Picking it is not a numerical detail --- it is the difference between a model with a unique equilibrium and one with a continuum.
\medskip

The next deck does this properly: log-linearization and undetermined coefficients, the matrix quadratic, the generalized Schur decomposition that solves it, and the eigenvalue count that decides the case.
\end{column}
\begin{column}{0.44\textwidth}
\begin{conceptbox}[title=\textbf{Next (5.A4)}]\small
Blanchard--Kahn counting; Klein's QZ solution; determinacy, indeterminacy and non-existence, measured on a New Keynesian model where the Taylor principle decides which one you get. And where Dynare fits.
\end{conceptbox}
\end{column}
\end{columns}
\end{frame}

\end{document}
